\subsection{معرفی و انگیزه}
این مکانیزم بر مبنای مبانی نظری پیشرفته در ادبیات توجه و فشرده‌سازی \cite{vaswani2017attention} توسعه یافته است.

مکانیزم توجه \lr{ViG Attention} با الهام از شبکه‌های عصبی گرافی بینایی (\lr{Vision GNN یا ViG}) برای غلبه بر پیچیدگی محاسباتی و حافظه‌ای مربع $\mathcal{O}(N^2)$ در پردازش توالی‌های طولانی طراحی شده است. در مدل‌های ترنسفورمر استاندارد، ضرب کامل ماتریس‌های پرس‌وجو و کلید نیازمند تشکیل نقشه توجه چگال با ابعاد $N \times N$ است که در توالی‌های طویل موجب انفجار ترافیک گذرگاه حافظه و مصرف شدید حافظه گرافیکی (\lr{VRAM}) می‌گردد.

ایده محوری در \lr{ViG Attention}، ترکیب دو شاخه محاسباتی مکمل از طریق یک ساختار همجوشی دروازه‌بندی‌شده (\lr{Gated Fusion}):
\begin{enumerate}
    \item \textbf{شاخه توجه گرافی توپولوژیکی (\lr{Topological Graph Attention}):} ایجاد یک گراف مجاورت خلوت روی توالی با پیوند دادن همسایگی‌های محلی (شعاع پنجره $\pm 16$) به همراه گره‌های لنگر سراسری (\lr{Global Anchor Tokens} شامل توکن ابتدایی $0$ و پایانی $N-1$) جهت تبادل مستقیم اطلاعات میان دو انتهای توالی.
    \item \textbf{شاخه توجه خطی سراسری (\lr{Global Linear Recurrence/Attention}):} بهره‌گیری از نگاشت‌های فعال‌سازی نامنفی کِرنل $\phi(x) = \operatorname{ELU}(x) + 1$ با ترتیب محاسباتی خطی و پایدارسازی مقیاس با مخرج مقید.
\end{enumerate}
این دو شاخه از طریق ضرایب تطبیقی پیش‌بینی‌شده توسط یک شبکه دروازه (\lr{Gate Network}) با یکدیگر ترکیب می‌شوند تا هم وابستگی‌های موضعی دقیق و هم الگوهای سراسری در سطح توالی حفظ گردند.
